\documentclass[a4paper,11pt]{article}

% ============================================================
%  Le Collé — Séance 1 : Déterminants, courbes paramétrées
%  © 2026 Ahmed Amine Ouqlani
% ============================================================

% ---------- Polices ----------
\usepackage[utf8]{inputenc}
\usepackage[T1]{fontenc}
\usepackage[sc]{mathpazo}          % texte et mathématiques : Palatino
\usepackage[scaled=0.88]{helvet}   % titres et repères : linéale
\linespread{1.07}
\usepackage{amsmath,amssymb,amsthm,mathrsfs,stmaryrd}
\usepackage{microtype}

% ---------- Mise en page ----------
\usepackage[a4paper,top=2.3cm,bottom=2.3cm,left=2.4cm,right=2.4cm,
            headheight=14pt,headsep=0.8cm,footskip=1.1cm]{geometry}
\setlength{\parindent}{0pt}
\setlength{\parskip}{0.45em plus 0.15em minus 0.1em}
\raggedbottom
\emergencystretch=1.5em
\clubpenalty=10000 \widowpenalty=10000 \displaywidowpenalty=10000 \brokenpenalty=10000

% ---------- Couleurs : un bleu de nuit, un ocre, des gris ----------
\usepackage{xcolor}
\definecolor{Ink}{HTML}{1C2533}
\definecolor{Bleu}{HTML}{1F3A5F}
\definecolor{BleuMoyen}{HTML}{4A6A8F}
\definecolor{BleuClair}{HTML}{9DB2C9}
\definecolor{Fond}{HTML}{F2F4F7}
\definecolor{Ocre}{HTML}{B07A26}
\definecolor{Brique}{HTML}{A23B2A}
\definecolor{Gris}{HTML}{5E6672}
\definecolor{Filet}{HTML}{C5CCD6}

% ---------- Packages ----------
\usepackage{enumitem}
\setlist{itemsep=0.25em,topsep=0.35em,parsep=0pt}
\setlist[itemize]{label=\textcolor{BleuMoyen}{\small\textbullet},leftmargin=1.4em}
\setlist[enumerate]{leftmargin=1.8em,label=\textsf{\small\color{Bleu}\arabic*.}}
\usepackage{booktabs,array,multicol}
\usepackage{needspace,float}
\usepackage[font=small,labelfont={sf,color=Bleu},labelsep=period,skip=6pt]{caption}
\usepackage{tikz}
\usetikzlibrary{arrows.meta,calc,decorations.markings}
\usepackage{pgfplots}
\pgfplotsset{compat=1.17}
\usepackage{tkz-tab}
\usepackage[most]{tcolorbox}
\usepackage{titlesec,titletoc}
\usepackage{fancyhdr,lastpage,etoolbox}
\usepackage[colorlinks=true,linkcolor=Bleu,urlcolor=Bleu,citecolor=Bleu,
            pdftitle={Le Collé — TD Séries numériques (énoncés)},pdfauthor={Ahmed Amine Ouqlani}]{hyperref}

\numberwithin{equation}{section}
\numberwithin{figure}{section}
\numberwithin{table}{section}

% ---------- Titres ----------
% Chapitres (sections) : grand numéro, titre en linéale, filet ocre
\titleformat{\section}[display]
  {\normalfont\sffamily\color{Ink}}
  {{\small\color{Gris}\textls[200]{\MakeUppercase{Chapitre}}}\\[-0.2em]%
   {\fontsize{64}{64}\selectfont\color{BleuClair}\thesection}}
  {0.4em}
  {\fontsize{25}{30}\selectfont\bfseries}
  [\vspace{0.5em}{\color{Ocre}\rule{3.2cm}{1.4pt}}]
\titlespacing*{\section}{0pt}{0pt}{1.6em}
\titleformat{name=\section,numberless}[display]
  {\normalfont\sffamily\color{Ink}}{}{0pt}{\fontsize{25}{30}\selectfont\bfseries}
  [\vspace{0.5em}{\color{Ocre}\rule{3.2cm}{1.4pt}}]
\titlespacing*{name=\section,numberless}{0pt}{0pt}{1.6em}

\titleformat{\subsection}
  {\needspace{5\baselineskip}\normalfont\sffamily\Large\bfseries\color{Ink}}
  {\color{Bleu}\thesubsection}{0.8em}{}
\titlespacing*{\subsection}{0pt}{1.8em}{0.5em}
\titleformat{\subsubsection}
  {\needspace{5\baselineskip}\normalfont\sffamily\large\bfseries\color{Ink}}
  {\color{Bleu}\thesubsubsection}{0.7em}{}
\titlespacing*{\subsubsection}{0pt}{1.6em}{0.4em}
\titleformat{\paragraph}[runin]{\normalfont\sffamily\bfseries\color{Ink}}{}{0pt}{}[.]
\titlespacing*{\paragraph}{0pt}{0.9em}{0.6em}
\setcounter{secnumdepth}{3}
\setcounter{tocdepth}{2}

% ---------- Table des matières ----------
\titlecontents{section}[0pt]{\addvspace{0.7em}\sffamily\bfseries\color{Ink}}
  {\makebox[2.2em][l]{\color{Bleu}\thecontentslabel}}{}{\hfill\contentspage}[\addvspace{0.2em}]
\titlecontents{subsection}[2.2em]{\small\linespread{1}\selectfont}
  {\makebox[2.6em][l]{\color{Gris}\thecontentslabel}}{}{\titlerule*[0.7pc]{.}\contentspage}

% ---------- En-têtes et pieds de page ----------
\pagestyle{fancy}
\fancyhf{}
\renewcommand{\sectionmark}[1]{\markboth{\thesection\quad#1}{}}
\fancyhead[L]{\footnotesize\sffamily\color{Gris}\leftmark}
\fancyhead[R]{\footnotesize\sffamily\scshape\color{Gris}Le Collé}
\fancyfoot[C]{\small\sffamily\color{Ink}\thepage}
\fancyfoot[L]{\scriptsize\sffamily\hyperlink{toc}{\color{Filet}sommaire}}
\renewcommand{\headrulewidth}{0.4pt}
\renewcommand{\headrule}{{\color{Filet}\hrule height\headrulewidth width\headwidth}}
\renewcommand{\footrulewidth}{0pt}
\fancypagestyle{plain}{\fancyhf{}\fancyfoot[C]{\small\sffamily\color{Ink}\thepage}\renewcommand{\headrulewidth}{0pt}}

% ---------- Énoncés ----------
\newtheoremstyle{enonce}{0pt}{0pt}{}{}{\sffamily\bfseries\color{Bleu}}{.}{0.5em}
  {\thmname{#1}\thmnumber{~#2}\thmnote{\ \normalfont\sffamily\color{Ink}(#3)}}
\newtheoremstyle{exemple}{0.9em}{0.9em}{}{}{\sffamily\bfseries\color{BleuMoyen}}{.}{0.5em}
  {\thmname{#1}\thmnumber{~#2}\thmnote{\ \normalfont\sffamily\color{Ink}(#3)}}
\newtheoremstyle{remarque}{0.7em}{0.7em}{}{}{\sffamily\itshape}{.}{0.5em}{\thmname{#1}\thmnote{\ (#3)}}
\theoremstyle{enonce}
\newtheorem{definition}{Définition}[section]
\newtheorem{theoreme}[definition]{Théorème}
\newtheorem{proposition}[definition]{Proposition}
\newtheorem{corollaire}[definition]{Corollaire}
\theoremstyle{exemple}
\newtheorem{exemple}[definition]{Exemple}
\theoremstyle{remarque}
\newtheorem*{remarque}{Remarque}

\tcbset{
  defi/.style={enhanced,breakable,colback=white,colframe=white,boxrule=0pt,arc=0pt,
     borderline west={1.2pt}{0pt}{Bleu},left=11pt,right=4pt,top=3pt,bottom=3pt,
     before skip=1.1em,after skip=1.1em},
  thm/.style={enhanced,breakable,colback=Fond,colframe=Fond,boxrule=0pt,arc=0pt,
     borderline west={2.4pt}{0pt}{Bleu},left=11pt,right=8pt,top=6pt,bottom=6pt,
     before skip=1.2em,after skip=1.2em}}
\tcolorboxenvironment{definition}{defi}
\tcolorboxenvironment{theoreme}{thm}
\tcolorboxenvironment{proposition}{thm}
\tcolorboxenvironment{corollaire}{thm}

\renewcommand{\proofname}{Démonstration}
\renewcommand{\qedsymbol}{\textcolor{BleuClair}{$\blacksquare$}}
\makeatletter
\renewenvironment{proof}[1][\proofname]{\par\pushQED{\qed}\normalfont
  \topsep0.5em\relax\trivlist\item[\hskip\labelsep\sffamily\itshape\color{Gris}#1.]\ignorespaces}
  {\popQED\endtrivlist\@endpefalse}
\makeatother

% ---------- Méthode : filets haut et bas, titre en bandeau ----------
\makeatletter\newcommand{\methlabelfix}{\protected@edef\@currentlabel{\thetcbcounter}}\makeatother
\newtcolorbox[auto counter,number within=section]{methode}[2][]{%
  enhanced,unbreakable,colback=white,colframe=white,boxrule=0pt,arc=0pt,
  borderline north={0.9pt}{0pt}{Bleu},borderline south={0.9pt}{0pt}{Bleu},
  left=4pt,right=4pt,top=8pt,bottom=8pt,before skip=1.4em,after skip=1.4em,
  before upper={\methlabelfix\phantomsection\addcontentsline{met}{metentry}{\protect\makebox[2.6em][l]{\thetcbcounter}#2}%
     {\sffamily\bfseries\small\color{white}\colorbox{Bleu}{\strut\ MÉTHODE~\thetcbcounter\ }}%
     \ \ {\sffamily\bfseries\color{Bleu}#2}\par\smallskip},#1}

% ---------- Repères pédagogiques : un filet de couleur, un mot-clé ----------
\newcommand{\definirbloc}[3]{% nom, couleur, libellé
  \newtcolorbox{#1}[1][]{enhanced,breakable,blanker,
    borderline west={1.4pt}{0pt}{#2},left=10pt,right=2pt,top=1pt,bottom=1pt,
    before skip=1em,after skip=1em,
    before upper={{\sffamily\bfseries\footnotesize\color{#2}\textls[120]{\MakeUppercase{#3}}}\quad},##1}}
\definirbloc{reflexe}{Bleu}{Réflexe}
\definirbloc{astuce}{Ocre}{Astuce}
\definirbloc{attention}{Brique}{Attention}
\definirbloc{enpratique}{BleuMoyen}{En pratique}
\definirbloc{plusloin}{Gris}{Pour aller plus loin}
\definirbloc{strategie}{Bleu}{Stratégie}
\definirbloc{lecture}{BleuMoyen}{Lire le schéma}
% compatibilité avec les parties reprises
\newenvironment{garde}{\begin{attention}}{\end{attention}}
\newenvironment{complementHP}{\begin{plusloin}}{\end{plusloin}}
\newenvironment{interpretation}{\begin{enpratique}}{\end{enpratique}}

% ---------- Exercices, corrigés ----------
\newcounter{exo}[section]
\renewcommand{\theexo}{\thesection.\arabic{exo}}
\newcommand{\niveau}[1]{{\color{Ocre}\footnotesize
  \ifcase#1\or$\bullet\circ\circ\circ\circ$\or$\bullet\bullet\circ\circ\circ$\or$\bullet\bullet\bullet\circ\circ$%
  \or$\bullet\bullet\bullet\bullet\circ$\or$\bullet\bullet\bullet\bullet\bullet$\fi}}
\newcommand{\serie}[1]{\par\needspace{8\baselineskip}\vspace{1.5em}%
  \noindent{\sffamily\bfseries\color{Bleu}#1}\par\nobreak\vspace{-0.4em}%
  \noindent{\color{Filet}\rule{\linewidth}{0.4pt}}\par\nobreak\vspace{0.1em}}
\newenvironment{exo}[3]{%
  \par\needspace{5\baselineskip}\vspace{0.75em}\refstepcounter{exo}\label{e:#1}%
  \def\exolab{#1}%
  \noindent{\sffamily\bfseries\color{Bleu}Exercice~\theexo}\quad\niveau{#2}%
  \hfill{\footnotesize\sffamily\color{Gris}#3}\par\nobreak\vspace{0.15em}}%
  {\par\nopagebreak\noindent\hfill{\scriptsize\sffamily\hyperref[c:\exolab]{\color{Gris}corrigé p.~\pageref{c:\exolab}}}\par}
\newenvironment{corrige}[1]{%
  \par\needspace{6\baselineskip}\vspace{1em}\phantomsection\label{c:#1}%
  \noindent{\sffamily\bfseries\color{Bleu}Exercice~\ref{e:#1}}\quad
  {\scriptsize\sffamily\hyperref[e:#1]{\color{Gris}énoncé p.~\pageref{e:#1}}}\par\nobreak\vspace{0.2em}}%
  {\hfill\textcolor{BleuClair}{$\square$}\par}
\newcommand{\etape}[1]{\par\smallskip\noindent{\sffamily\bfseries\small\color{BleuMoyen}#1.}\ }
\newenvironment{solution}{\par\noindent{\sffamily\bfseries\small\color{BleuMoyen}Solution.}\ }{\hfill\textcolor{BleuClair}{$\square$}\par}
\newcommand{\chercher}{\par\smallskip\noindent{\color{Filet}\rule[0.5ex]{1.2cm}{0.4pt}}\
  {\footnotesize\sffamily\itshape\color{Gris}laisser chercher l'étudiant avant de poursuivre}\par\smallskip}

% ---------- Guide de séance ----------

\newcommand{\marge}[1]{\leavevmode\marginpar{\raggedleft\scriptsize\sffamily\color{Gris}\textls[100]{\MakeUppercase{#1}}}\ignorespaces}
\newcommand{\Etape}[4]{%
  \subsection[Étape #1 --- #3 {\color{Gris}(#2)}]{#3}\label{etape:#1}%
  \noindent{\small\sffamily\color{Gris}{\color{Ocre}\bfseries Étape #1}\quad·\quad #2\quad·\quad #4}\par\medskip}
\newcommand{\question}[1]{\par\smallskip\noindent{\sffamily\bfseries\small\color{BleuMoyen}Question.}\ \textit{#1}\par\smallskip}
\newcommand{\rappel}[1]{\par\smallskip\noindent\phantomsection\addcontentsline{rap}{rapentry}{#1}%
  {\sffamily\bfseries\scriptsize\colorbox{Fond}{\color{Gris}\strut RAPPEL}}\ {\sffamily\small\color{Gris}#1.}\ }
\newcommand{\pourquoi}{\par\noindent{\sffamily\bfseries\small\color{Ocre}Pourquoi ?}\ }

% Ouverture de chapitre : \ouverture{objectifs}{prérequis}{durée}
\newcommand{\ouverture}[3]{%
  \begin{tcolorbox}[enhanced,colback=Fond,colframe=Fond,arc=0pt,boxrule=0pt,left=12pt,right=12pt,top=10pt,bottom=10pt,
     before skip=1.2em,after skip=1.6em]
  \begin{minipage}[t]{0.56\linewidth}
  {\sffamily\bfseries\small\color{Bleu}OBJECTIFS}\par\vspace{0.3em}\small #1
  \end{minipage}\hfill
  \begin{minipage}[t]{0.38\linewidth}
  {\sffamily\bfseries\small\color{Bleu}PRÉREQUIS}\par\vspace{0.3em}\small #2
  \ifx\relax#3\relax\else\par\vspace{0.6em}{\sffamily\bfseries\small\color{Bleu}DURÉE}\par\vspace{0.2em}\small #3\fi
  \end{minipage}
  \end{tcolorbox}}
% Fin de chapitre
\newenvironment{essentiel}{\par\needspace{12\baselineskip}\vspace{1.5em}%
  \begin{tcolorbox}[enhanced,unbreakable,colback=white,colframe=white,boxrule=0pt,arc=0pt,
   borderline north={1.4pt}{0pt}{Ocre},left=2pt,right=2pt,top=8pt,bottom=4pt]
  {\sffamily\bfseries\color{Ocre}L'ESSENTIEL DU CHAPITRE}\par\medskip}{\end{tcolorbox}}

% ---------- Notations ----------
\newcommand{\K}{\mathbb{K}}
\newcommand{\R}{\mathbb{R}}
\newcommand{\C}{\mathbb{C}}
\newcommand{\N}{\mathbb{N}}
\newcommand{\Z}{\mathbb{Z}}
\newcommand{\Mn}[1]{\mathcal{M}_{n}(#1)}
\newcommand{\GLn}[1]{\mathrm{GL}_n(#1)}
\newcommand{\LE}{\mathcal{L}(E)}
\newcommand{\rg}{\operatorname{rg}}
\newcommand{\vect}{\operatorname{Vect}}
\newcommand{\Mat}{\operatorname{Mat}}
\newcommand{\Id}{\operatorname{Id}}
\newcommand{\sh}{\operatorname{sh}}
\newcommand{\ch}{\operatorname{ch}}
\renewcommand{\th}{\operatorname{th}}
\DeclareMathOperator{\argsh}{argsh}
\newcommand{\vv}[1]{\overrightarrow{#1}}
\newcommand{\ve}[1]{\vec{#1}}
\newcommand{\ient}[2]{\llbracket #1,#2\rrbracket}
\newcommand{\dt}[1]{\left|\begin{matrix}#1\end{matrix}\right|}
\newcommand{\dis}{\displaystyle}

% ---------- Figures ----------
\pgfplotsset{trig format plots=rad}
\pgfplotsset{courbe/.style={axis lines=middle,axis equal image,
  axis line style={Gris,-{Stealth[length=5pt]}},
  xlabel={$x$},ylabel={$y$},ticklabel style={font=\scriptsize\sffamily,color=Gris},
  every axis x label/.style={at={(ticklabel* cs:1)},anchor=west,font=\small},
  every axis y label/.style={at={(ticklabel* cs:1)},anchor=south,font=\small},
  samples=300,smooth,width=0.6\linewidth}}
\tikzset{arcetude/.style={Bleu,line width=1.25pt},
         reste/.style={BleuClair,line width=1pt},
         tgte/.style={Ocre,line width=0.9pt},
         asy/.style={Gris,dashed,line width=0.6pt},
         pt/.style={circle,fill=Ink,inner sep=1.2pt}}

% ---------- Navigation ----------
\makeatletter
\newcommand*\l@rapentry{\@dottedtocline{1}{0em}{0em}}
\newcommand*\l@metentry{\@dottedtocline{1}{0em}{0em}}
\newcommand{\listemethodes}{\@starttoc{met}}
\newcommand{\listerappels}{\@starttoc{rap}}
\makeatother
% onglets de chapitre sur le bord droit
\usepackage{eso-pic}
\newif\ifonglets \ongletsfalse
\newcommand{\dessineronglet}{%
  \edef\ongcol{\ifnum\value{section}=3 Ocre\else Bleu\fi}%
  \edef\ongnum{\the\value{section}}%
  \pgfmathsetmacro{\ongy}{3.2+2.3*(\ongnum-1)}%
  \begin{tikzpicture}[remember picture,overlay]
  \fill[\ongcol] ([xshift=-0.75cm,yshift=-\ongy cm]current page.north east)
     rectangle ([yshift=-\ongy cm-1.9cm]current page.north east);
  \node[text=white,font=\sffamily\bfseries\large] at ([xshift=-0.37cm,yshift=-\ongy cm-0.95cm]current page.north east)
     {\ongnum};
  \end{tikzpicture}}
\AddToShipoutPictureBG{\ifonglets\ifnum\value{section}>0 \dessineronglet\fi\fi}
% blocs de navigation et d'enseignement
\definirbloc{commencer}{Gris}{Si vous commencez ici}
\definirbloc{horsprogramme}{Brique}{Hors programme}
\newcommand{\HP}{\,{\sffamily\bfseries\scriptsize\textcolor{Brique}{[HP]}}}
\newcommand{\source}[1]{{\footnotesize\sffamily\color{Gris}\,(#1)}}
\newcommand{\coupdepouce}[1]{\par\vspace{-0.2em}\noindent{\footnotesize\sffamily\color{Gris}%
  \textbf{Coup de pouce} (à donner si l'étudiant bloque) : #1}\par}
\newcommand{\pointarret}[1]{\par\vspace{1em}\noindent{\color{Filet}\rule[0.5ex]{\linewidth}{0.4pt}}\par\nobreak
  \noindent{\footnotesize\sffamily\color{Ocre}\bfseries POINT D'ARRÊT POSSIBLE}\quad{\footnotesize\sffamily\color{Gris}#1}\par\vspace{0.6em}}
\newcommand{\demarche}[1]{\noindent{\sffamily\bfseries\small\color{BleuMoyen}Démarche.}\ \textit{#1}\par\smallskip}

% ---------- Corrigés : idées, astuces, pièges, à retenir ----------
\newcommand{\choix}[1]{\par\noindent{\sffamily\bfseries\small\color{BleuMoyen}Choix de la méthode.}\ #1\par\smallskip}
\newcommand{\gain}[1]{\par\noindent{\sffamily\bfseries\small\color{Ocre}Astuce.}\ #1\par}
\newcommand{\piege}[1]{\par\noindent{\sffamily\bfseries\small\color{Brique}Piège.}\ #1\par}
\newtcolorbox{retenirbox}{enhanced,breakable,colback=Fond,colframe=Fond,arc=0pt,boxrule=0pt,
  borderline west={2pt}{0pt}{Ocre},left=8pt,right=6pt,top=4pt,bottom=4pt,before skip=0.7em,after skip=0.5em}
\newcommand{\retenir}[1]{\begin{retenirbox}{\sffamily\bfseries\footnotesize\color{Ocre}À RETENIR}\quad #1\end{retenirbox}}
\newcommand{\TD}[1]{{\color{Ocre}\bfseries TD\,#1}\ ·\ }
\newcommand{\q}[1]{\par\smallskip\noindent{\sffamily\bfseries\color{Bleu}#1}\ }
\providecommand{\iddots}{\mathinner{\mkern1mu\raise1pt\hbox{.}\mkern2mu\raise4pt\hbox{.}\mkern2mu\raise7pt\hbox{.}\mkern1mu}}

\providecommand{\Ker}{\operatorname{Ker}}\renewcommand{\Im}{\operatorname{Im}}\providecommand{\tr}{\operatorname{tr}}\providecommand{\Sp}{\operatorname{Sp}}

% ================================================================ TD : mise en page propre au livret d'exercices
\fancyhf{}
\fancyhead[L]{\footnotesize\sffamily\color{Gris}TD · Séries numériques · énoncés}
\fancyhead[R]{\footnotesize\sffamily\scshape\color{Gris}Le Collé}
\fancyfoot[C]{\small\sffamily\color{Ink}\thepage\,/\,\pageref{LastPage}}
\fancypagestyle{plain}{\fancyhf{}\fancyfoot[C]{\small\sffamily\color{Ink}\thepage\,/\,\pageref{LastPage}}\renewcommand{\headrulewidth}{0pt}}
\makeatletter
\newcommand{\partie}[3]{% numéro, titre, description
  \par\needspace{10\baselineskip}\vspace{1.4em}\phantomsection\label{partie:#1}\pdfbookmark[0]{Partie #1}{partie.#1}%
  {\sffamily\small\color{Gris}\textls[200]{PARTIE #1}}\par\nobreak\vspace{0.1em}
  {\sffamily\bfseries\LARGE\color{Ink}#2}\par\nobreak\vspace{0.3em}
  {\color{Ocre}\rule{2.6cm}{1.2pt}}\par\nobreak\vspace{0.2em}
  {\small\color{Gris}#3}\par\nobreak\vspace{0.6em}}
\newcommand{\exercice}[7]{% numéro, difficulté, titre, origine, notions, identifiant, page du corrigé
  \par\needspace{8\baselineskip}\vspace{1.1em}\phantomsection
  \def\@currentlabel{#1}\label{e:#6}\pdfbookmark[1]{Exercice #1}{e.#6}%
  \noindent{\sffamily\bfseries\large\color{Bleu}Exercice~#1}\quad\niveau{#2}\hfill{\footnotesize\sffamily\color{Gris}#4}\par\nobreak
  \vspace{0.1em}\noindent{\sffamily\bfseries\color{Ink}#3}\par\nobreak
  \vspace{0.1em}\noindent{\footnotesize\sffamily\color{Gris}Notions : #5.\hfill\mbox{corrigé p.~#7}\par}\nobreak
  \vspace{0.35em}}
\newcommand{\renvoi}[1]{\par\nopagebreak\vspace{0.2em}\noindent\hfill{\scriptsize\sffamily\color{Gris}#1}\par}
\newcommand{\place}{\par\vspace{1.6cm plus 0.6cm minus 0.6cm}}
\newcommand{\correction}[4]{% numéro, titre, identifiant, renvoi
  \par\needspace{7\baselineskip}\vspace{1.3em}\phantomsection
  \def\@currentlabel{#1}\label{e:#3}\pdfbookmark[1]{Exercice #1}{c.#3}%
  {\color{Filet}\hrule height 0.4pt}\vspace{0.5em}
  \noindent{\sffamily\bfseries\large\color{Bleu}Exercice~#1}\quad{\sffamily\bfseries\color{Ink}#2}\hfill{\scriptsize\sffamily\color{Gris}#4}\par\nobreak
  \vspace{0.4em}}
\newcommand{\sommaireligne}[3]{\noindent\hbox to 2.2em{\hfill\sffamily\color{Bleu}#1.}\enspace\begin{minipage}[t]{0.8\linewidth}\raggedright#2\end{minipage}\hfill\pageref{e:#3}\par}
\makeatother
\setlength{\parskip}{0.5em plus 0.15em minus 0.1em}

\begin{document}
\thispagestyle{plain}
{\sffamily\small\color{Gris}\textls[150]{LE COLLÉ · MATHÉMATIQUES PT · CHAPITRE 7}}\par\vspace{0.2em}
{\sffamily\bfseries\fontsize{26}{30}\selectfont\color{Ink}TD --- Séries numériques}\par\vspace{0.3em}
{\sffamily\Large\color{BleuMoyen}Énoncés}\par\vspace{0.4em}
{\color{Ocre}\rule{3.2cm}{1.4pt}}\par\vspace{0.8em}
{\small Ce livret rassemble 25 exercices, classés en quatre parties, du plus direct au plus exigeant.
Il reprend \textbf{tous les exercices du livret} (proposés à la fin de chaque partie du cours), complétés par des
\textbf{exercices complémentaires}, absents du livret. L'origine de chaque exercice est indiquée en haut à droite :
\emph{« Livret, exercice~1.5 »} ou \emph{« Exercice complémentaire »}. La difficulté va de \niveau{1} à \niveau{5}.
Les corrections sont dans un livret séparé, \texttt{TD\_series\_corrections.pdf} ; la page de chaque correction est indiquée en tête de l'énoncé.}

\vspace{0.6em}{\sffamily\bfseries\color{Bleu}Sommaire}\par\vspace{0.2em}{\small\setlength{\parskip}{1.5pt}\setlength{\columnsep}{1.2cm}\begin{multicols}{2}\raggedcolumns
\vspace{0.5em}\noindent\begin{minipage}[t]{0.86\linewidth}\raggedright\leavevmode\sffamily\bfseries\color{Ink}I. Applications directes du cours\end{minipage}\hfill\pageref{partie:I}\par\vspace{0.15em}
\sommaireligne{1}{Suites classiques}{sn-1}
\sommaireligne{2}{Équivalents}{sn-2}
\sommaireligne{3}{Théorème des gendarmes}{sn-3}
\sommaireligne{4}{Un développement asymptotique de $(1+\frac1n)^n$}{s-equiv}
\sommaireligne{5}{Télescopages}{s-tele}
\sommaireligne{6}{Somme par éléments simples}{sn-5}
\sommaireligne{7}{Nature de séries à termes positifs}{s-nature}
\sommaireligne{8}{Nature : comparaison, équivalents, d'Alembert}{sn-4}
\vspace{0.5em}\noindent\begin{minipage}[t]{0.86\linewidth}\raggedright\leavevmode\sffamily\bfseries\color{Ink}II. Exercices classiques et méthodes incontournables\end{minipage}\hfill\pageref{partie:II}\par\vspace{0.15em}
\sommaireligne{9}{Constante d'Euler par suites adjacentes}{s-gamma}
\sommaireligne{10}{Une suite étudiée par une série}{s-gamma2}
\sommaireligne{11}{Règle de d'Alembert}{s-dal}
\sommaireligne{12}{Série géométrique dérivée}{sn-6}
\sommaireligne{13}{Séries de Bertrand}{sn-8}
\sommaireligne{14}{La série harmonique alternée}{s-alt}
\sommaireligne{15}{Une série semi-convergente}{sn-9}
\sommaireligne{16}{Produit de Cauchy de séries géométriques}{s-cauchy}
\vspace{0.5em}\noindent\begin{minipage}[t]{0.86\linewidth}\raggedright\leavevmode\sffamily\bfseries\color{Ink}III. Exercices de synthèse et approfondissement\end{minipage}\hfill\pageref{partie:III}\par\vspace{0.15em}
\sommaireligne{17}{Équivalent du reste d'une série de Riemann}{s-reste}
\sommaireligne{18}{Développements asymptotiques}{s-asympt}
\sommaireligne{19}{La famille $\sum\frac1{n^\alpha\ln n}$}{sn-16}
\sommaireligne{20}{Somme d'une série alternée par une intégrale}{sn-7}
\sommaireligne{21}{Une série définie par récurrence}{sn-11}
\sommaireligne{22}{Une série trigonométrique surprenante}{sn-12}
\vspace{0.5em}\noindent\begin{minipage}[t]{0.86\linewidth}\raggedright\leavevmode\sffamily\bfseries\color{Ink}IV. Exercices difficiles et type concours\end{minipage}\hfill\pageref{partie:IV}\par\vspace{0.15em}
\sommaireligne{23}{Reste d'une série alternée}{sn-13}
\sommaireligne{24}{Un produit de Cauchy qui diverge}{sn-14}
\sommaireligne{25}{La suite $u_{n+1}=\sin u_n$}{sn-15}
\end{multicols}}
\clearpage
\partie{I}{Applications directes du cours}{Suites classiques, équivalents, gendarmes, télescopages, nature de séries à termes positifs.}
\exercice{1}{1}{Suites classiques}{Exercice complémentaire}{arithmético-géométrique, récurrence linéaire d'ordre 2}{sn-1}{1}
Déterminer le terme général des suites définies par :
\begin{enumerate}
\item $u_0=0$ et $u_{n+1}=-u_n+2$ ;
\item $u_0=1$, $u_1=2$ et $u_{n+2}+u_{n+1}+u_n=0$ ;
\item $u_0=0$, $u_1=1$ et $u_{n+2}-3u_{n+1}+2u_n=0$.
\end{enumerate}
\place
\exercice{2}{1}{Équivalents}{Exercice complémentaire}{équivalents, quantité conjuguée, développement limité}{sn-2}{1}
Donner un équivalent simple de :
\[\text{(a)}\ \sqrt{n^2+n+1}-n\qquad\text{(b)}\ \sin\frac{n\pi}{n+1}\qquad\text{(c)}\ \sqrt{n+(-1)^n}-\sqrt n\qquad\text{(d)}\ \ln(n+1)-\ln n.\]
\place
\exercice{3}{2}{Théorème des gendarmes}{Exercice complémentaire}{encadrement, partie entière}{sn-3}{2}
Étudier la limite de $u_n=\sum_{k=1}^n\frac1{\sqrt{n^2+k}}$, puis, pour $x\in\R$ fixé, de $v_n=\frac1{n^2}\sum_{k=1}^n\lfloor kx\rfloor$.
\place
\exercice{4}{2}{Un développement asymptotique de $(1+\frac1n)^n$}{Livret, exercice 1.1}{équivalents, développement limité, exponentielle}{s-equiv}{2}
Montrer que $\left(1+\frac1n\right)^n\to e$, puis que $\left(1+\frac1n\right)^n-e\sim-\frac e{2n}$.
\place
\exercice{5}{2}{Télescopages}{Livret, exercice 1.3}{série télescopique, éléments simples, logarithme}{s-tele}{2}
\begin{enumerate}
\item Décomposer $\frac1{n(n+1)(n+2)}$ en éléments simples, puis calculer $\sum_{n=1}^{+\infty}\frac1{n(n+1)(n+2)}$.
\item Montrer que $\sum\ln\left(1+\frac1n\right)$ diverge.
\item Montrer que $\sum_{n\geq2}\ln\left(1-\frac1{n^2}\right)$ converge et calculer sa somme.
\end{enumerate}
\place
\exercice{6}{1}{Somme par éléments simples}{Exercice complémentaire}{éléments simples, série télescopique}{sn-5}{2}
Justifier la convergence et calculer $\sum_{n=1}^{+\infty}\frac1{n(n+2)}$.
\place
\exercice{7}{2}{Nature de séries à termes positifs}{Livret, exercice 1.5}{équivalents, Riemann, croissances comparées}{s-nature}{3}
Déterminer la nature des séries de terme général :
\[\text{(a)}\ \frac{n^2+1}{n^4+3}\qquad\text{(b)}\ \frac1{\sqrt{n(n+1)}}\qquad\text{(c)}\ 1-\cos\frac1n\qquad\text{(d)}\ \frac1{\ln n+e^n}\qquad\text{(e)}\ n^2e^{-\sqrt n}\qquad\text{(f)}\ \frac{\ln n}n.\]
\place
\exercice{8}{2}{Nature : comparaison, équivalents, d'Alembert}{Exercice complémentaire}{d'Alembert, équivalents, divergence grossière}{sn-4}{3}
Déterminer la nature des séries de terme général :
\[\text{(a)}\ \frac{2^nn!}{n^n}\quad\text{(b)}\ \frac{n^4}{3^n}\quad\text{(c)}\ \sin\frac1n+\ln\left(1-\frac1n\right)\quad\text{(d)}\ n\sin\frac1n\quad
\text{(e)}\ \sin(e^{-3n})\quad\text{(f)}\ \frac1{\sqrt[3]{n^2-1}}-\frac1{\sqrt[3]{n^2+1}}\quad\text{(g)}\ \frac1{\sqrt n\cos^2n}.\]
\place
\clearpage
\partie{II}{Exercices classiques et méthodes incontournables}{Constante d'Euler, d'Alembert, séries géométriques dérivées, comparaison série-intégrale, séries alternées, produit de Cauchy.}
\exercice{9}{2}{Constante d'Euler par suites adjacentes}{Livret, exercice 1.2}{suites adjacentes, série harmonique}{s-gamma}{3}
On pose $H_n=\sum_{k=1}^n\frac1k$, $u_n=H_n-\ln n$ et $v_n=u_n-\frac1n$ ($n\geq1$). En utilisant $\ln(1+x)\leq x$ pour $x>-1$,
montrer que $(u_n)$ et $(v_n)$ sont adjacentes. On note $\gamma$ leur limite commune : montrer que $\gamma>0$, et que
$H_n=\ln n+\gamma+o(1)$.
\place
\exercice{10}{2}{Une suite étudiée par une série}{Livret, exercice 1.4}{série télescopique, équivalent, Riemann}{s-gamma2}{4}
On reprend $u_n=H_n-\ln n$ (exercice~\ref{e:s-gamma}). Donner un équivalent de $u_n-u_{n-1}$ ; en déduire, sans suites
adjacentes, que $(u_n)$ converge.
\place
\exercice{11}{2}{Règle de d'Alembert}{Livret, exercice 1.6}{d'Alembert, factorielles, cas douteux}{s-dal}{4}
\begin{enumerate}
\item Soit $a>0$. Montrer que $\sum\frac{a^n}{n!}$ converge ; en déduire que $a^n=o(n!)$.
\item Nature de $\sum\frac{(2n)!}{(n!)^2\,4^n}$ ? Pourquoi d'Alembert ne permet-il pas de conclure ? (On admettra que
$\frac{(2n)!}{(n!)^24^n}\sim\frac1{\sqrt{\pi n}}$.)
\end{enumerate}
\place
\exercice{12}{2}{Série géométrique dérivée}{Exercice complémentaire}{série géométrique, dérivation, linéarité}{sn-6}{4}
\begin{enumerate}
\item Pour $x\neq1$, calculer $\sum_{k=0}^nx^k$, puis $\sum_{k=1}^nkx^{k-1}$ en dérivant.
\item En déduire que, pour $|x|<1$, $\sum_{k=1}^{+\infty}kx^{k-1}=\frac1{(1-x)^2}$.
\item Justifier la convergence et calculer $\sum_{k=0}^{+\infty}\big(k-(-1)^k\big)3^{-k}$.
\end{enumerate}
\place
\exercice{13}{2}{Séries de Bertrand}{Exercice complémentaire}{comparaison série-intégrale, logarithme}{sn-8}{4}
\begin{enumerate}
\item Par comparaison avec une intégrale, montrer que $\sum_{n\geq2}\frac1{n(\ln n)^2}$ converge.
\item Montrer que $\sum_{k=2}^n\frac1{k\ln k}\sim\ln(\ln n)$.
\end{enumerate}
\place
\exercice{14}{2}{La série harmonique alternée}{Livret, exercice 1.8}{séries alternées, majoration du reste, intégration}{s-alt}{5}
\begin{enumerate}
\item Montrer que $\sum_{n\geq1}\frac{(-1)^{n+1}}n$ converge. En intégrant $\sum_{k=0}^{n-1}(-t)^k=\frac{1-(-t)^n}{1+t}$ sur $[0,1]$,
montrer que sa somme vaut $\ln2$.
\item Combien de termes faut-il sommer pour approcher $\ln 2$ à $10^{-3}$ près par cette série ?
\end{enumerate}
\place
\exercice{15}{2}{Une série semi-convergente}{Exercice complémentaire}{semi-convergence, comparaison série-intégrale, séries alternées}{sn-9}{5}
\begin{enumerate}
\item Montrer que $\sum_{k=1}^n\frac{\ln k}k\sim\frac{(\ln n)^2}2$.
\item Montrer que $\sum\frac{(-1)^n\ln n}n$ converge, mais pas absolument.
\end{enumerate}
\place
\exercice{16}{2}{Produit de Cauchy de séries géométriques}{Livret, exercice 1.10}{produit de Cauchy, série géométrique}{s-cauchy}{5}
\begin{enumerate}
\item Calculer $\left(\sum_{n=0}^{+\infty}\frac1{2^n}\right)^2$ et $\sum_{n=0}^{+\infty}\frac1{4^n}$ : que constate-t-on ?
\item Pour $|x|<1$, calculer $\sum_{n=0}^{+\infty}nx^n$, puis $\sum_{n=0}^{+\infty}\frac{n}{2^n}$.
\end{enumerate}
\place
\clearpage
\partie{III}{Exercices de synthèse et approfondissement}{Équivalent d'un reste, développements asymptotiques, familles de séries, séries définies par récurrence.}
\exercice{17}{3}{Équivalent du reste d'une série de Riemann}{Livret, exercice 1.7}{comparaison série-intégrale, reste}{s-reste}{6}
Soit $\alpha>1$. Montrer, par comparaison avec une intégrale, que $\displaystyle\sum_{k=n+1}^{+\infty}\frac1{k^\alpha}\sim\frac1{(\alpha-1)n^{\alpha-1}}$.
\place
\exercice{18}{3}{Développements asymptotiques}{Livret, exercice 1.9}{développement asymptotique, séries alternées, semi-convergence}{s-asympt}{6}
Déterminer la nature des séries de terme général
\[\text{(a)}\ \ln\left(1+\frac{(-1)^n}{\sqrt n}\right)\ (n\geq2)\qquad\text{(b)}\ \frac{(-1)^n}{n+(-1)^n}\ (n\geq2)\qquad\text{(c)}\ \frac{(-1)^n}{n^\alpha}\ (\alpha\in\R).\]
\place
\exercice{19}{2}{La famille $\sum\frac1{n^\alpha\ln n}$}{Exercice complémentaire}{Riemann, croissances comparées, série-intégrale}{sn-16}{6}
Selon $\alpha\in\R$, déterminer la nature de $\sum_{n\geq2}\frac1{n^\alpha\ln n}$.
\place
\exercice{20}{3}{Somme d'une série alternée par une intégrale}{Exercice complémentaire}{séries alternées, intégration, éléments simples}{sn-7}{6}
\begin{enumerate}
\item Montrer que $\sum\frac{(-1)^n}{3n+1}$ converge.
\item Montrer que, pour tout $n$, $\displaystyle\sum_{k=0}^n\frac{(-1)^k}{3k+1}=\int_0^1\frac{dt}{1+t^3}-\int_0^1\frac{(-t^3)^{n+1}}{1+t^3}\,dt$.
\item En déduire $\displaystyle\sum_{n=0}^{+\infty}\frac{(-1)^n}{3n+1}=\int_0^1\frac{dt}{1+t^3}$, puis calculer cette intégrale
(décomposer $\frac1{1+t^3}=\frac\alpha{1+t}+\frac{\beta t+\gamma}{t^2-t+1}$).
\end{enumerate}
\place
\exercice{21}{3}{Une série définie par récurrence}{Exercice complémentaire}{suite récurrente, série télescopique, équivalents}{sn-11}{7}
Soit $u_0\in\,]-1,1[$ et $u_{n+1}=\frac{1+u_n^2}2$. On pose $w_n=1-u_n$.
\begin{enumerate}
\item Montrer que $u_n\in[0,1[$ pour $n\geq1$, que $(u_n)$ est croissante, et qu'elle tend vers $1$.
\item Montrer que $w_{n+1}=w_n-\frac{w_n^2}2$. En déduire que $\sum w_n^2$ converge.
\item Déterminer la nature de $\sum\ln\frac{w_{n+1}}{w_n}$, puis celle de $\sum w_n$.
\end{enumerate}
\place
\exercice{22}{3}{Une série trigonométrique surprenante}{Exercice complémentaire}{binôme, série géométrique, équivalent}{sn-12}{7}
\begin{enumerate}
\item Montrer que la série de terme général $u_n=\sin\big((2-\sqrt3)^n\pi\big)$ est à termes positifs et convergente.
\item On pose $A_n=(2-\sqrt3)^n+(2+\sqrt3)^n$. Montrer que $A_n$ est un entier pair (développer par la formule du binôme).
\item En déduire la nature de la série de terme général $v_n=\sin\big((2+\sqrt3)^n\pi\big)$.
\end{enumerate}
\place
\clearpage
\partie{IV}{Exercices difficiles et type concours}{Reste d'une série alternée, produit de Cauchy divergent, la suite $u_{n+1}=\sin u_n$.}
\exercice{23}{4}{Reste d'une série alternée}{Exercice complémentaire}{reste, séries alternées, développement asymptotique}{sn-13}{7}
Soit $R_n=\sum_{k=n+1}^{+\infty}\frac{(-1)^k}k$ le reste de la série $\sum_{k\geq1}\frac{(-1)^k}k$.
\begin{enumerate}
\item Calculer $R_n-R_{n+1}$. Montrer que $R_n+R_{n+1}=\sum_{k=n+1}^{+\infty}\frac{(-1)^k}{k(k+1)}$ et que $|R_n+R_{n+1}|\leq\frac1{(n+1)(n+2)}$.
\item En déduire que $R_n\sim\frac{(-1)^{n+1}}{2n}$, puis que la série $\sum R_n$ converge.
\end{enumerate}
\place
\exercice{24}{4}{Un produit de Cauchy qui diverge}{Exercice complémentaire}{produit de Cauchy, semi-convergence}{sn-14}{8}
Soit $u_n=\frac{(-1)^n}{\sqrt{n+1}}$. Montrer que $\sum u_n$ converge. On note $\sum w_n$ le produit de Cauchy de $\sum u_n$ par elle-même.
Montrer que $|w_n|\geq1$ pour tout $n$ (on pourra majorer $(p+1)(n-p+1)$). Conclure.
\place
\exercice{25}{4}{La suite $u_{n+1}=\sin u_n$}{Exercice complémentaire}{Cesàro, équivalent, suite récurrente}{sn-15}{8}
Soit $u_0\in\,]0,\pi[$ et $u_{n+1}=\sin u_n$.
\begin{enumerate}
\item Montrer que $u_n\in\,]0,\pi[$, que $(u_n)$ décroît et tend vers $0$.
\item \emph{Lemme de Cesàro.} Montrer que si $x_n\to\ell\in\R$, alors $\frac1n\sum_{k=1}^nx_k\to\ell$.
\item Montrer que $\frac1{u_{n+1}^2}-\frac1{u_n^2}\to\frac13$. En déduire que $u_n\sim\sqrt{\frac3n}$.
\item Nature de $\sum u_n$, $\sum u_n^2$, $\sum u_n^3$ ?
\end{enumerate}
\place
\end{document}
